% TeX Live's cmtiup package identifies itself as ``cmtiup.sty'', while
% aomart requests ``cmtiup''.  Register it here and install its OT1 font
% shapes below to avoid the resulting spurious package-name warning.
\makeatletter
\expandafter\def\csname ver@cmtiup.sty\endcsname
  {2016/02/14 v2.1 Upright digits and punctuation}
\makeatother
\documentclass{aomart}

\usepackage{amsmath,amssymb,mathtools,microtype}

\DeclareFontShape{OT1}{cmr}{m}{it}{%
  <5><6><7>cmtiup7<8>cmtiup8<9>cmtiup9%
  <10><10.95>cmtiup10<12><14.4><17.28><20.74><24.88>cmtiup12}{}
\DeclareFontShape{OT1}{cmr}{bx}{it}{%
  <5><6><7><8><9><10><10.95><12><14.4><17.28><20.74><24.88>%
  cmbxtiup10}{}

\allowdisplaybreaks[2]
\setlength{\headheight}{20pt}
\setlength{\emergencystretch}{3em}
\raggedbottom
\makeatletter
\pdfstringdefDisableCommands{\let\unskip\@empty}
\makeatother

% Production identifiers are assigned by the Annals after acceptance.
% Suppress the class defaults in the submission PDF.
\doinumber{}
\copyrightnote{}
\fancypagestyle{firstpage}{%
  \fancyhf{}%
  \cfoot{\footnotesize\thepage}}

\newtheorem{theorem}{Theorem}[section]
\newtheorem{lemma}[theorem]{Lemma}
\newtheorem{proposition}[theorem]{Proposition}
\newtheorem{corollary}[theorem]{Corollary}

\theoremstyle{definition}
\newtheorem{remark}[theorem]{Remark}

\newcommand{\C}{\mathbb C}
\newcommand{\D}{\mathbb D}
\newcommand{\Torus}{\mathbb T}
\newcommand{\Mn}{\mathbb M_n(\C)}
\newcommand{\alg}{\operatorname{alg}}
\newcommand{\dist}{\operatorname{dist}}
\newcommand{\RePart}{\operatorname{Re}}
\newcommand{\diag}{\operatorname{diag}}

\title[The numerical range is a 2-spectral set]
  {The numerical range is a 2-spectral set}
\author{Shanmu Jin}
\address{Department of Neurosurgery, Peking Union Medical College
  Hospital, No.~1 Shuaifuyuan, Dongcheng District, Beijing 100730,
  China}
\email{jinshanmu1129@pku.edu.cn}

\keyword{Crouzeix conjecture}
\keyword{numerical range}
\keyword{spectral set}
\keyword{matrix-valued Herglotz kernel}
\keyword{double-layer potential}
\subject{primary}{matsc2020}{47A25}
\subject{secondary}{matsc2020}{47A12}
\subject{secondary}{matsc2020}{15A60}

\begin{document}

\begin{abstract}
Let $A$ be a complex square matrix and let $W(A)$ be its numerical
range.  We prove that
\[
  \lVert p(A)\rVert\leq
  2\max_{z\in W(A)}\lvert p(z)\rvert
\]
for every complex polynomial $p$, and hence for every rational function
with no pole on $W(A)$.  The constant is optimal.  The proof combines
the classical positive double-layer functional calculus with a new
positive-real completion principle.  The argument first treats
algebra-generating, simple-spectrum approximants.  For these auxiliary
objects, the double-layer map applied to the Cayley family completes the
resolvent of $T=f(B)$ by a term in the algebra generated by $T^*$.
Positivity of the resulting matrix-valued Herglotz kernel is sampled at
points determined by the eigenvalues of $T$, and a compensating sample
at the origin cancels the completion exactly.  The remaining
Pick-matrix inequality compares two weighted observability Gramians
and, through a Stein identity, gives $\lVert T\rVert\leq2$.
Perturbation, density, and smooth convex outer approximation then pass
the estimate to arbitrary matrices and their exact numerical ranges.
Thus no simple-spectrum, diagonalizability, algebra-generation, or
boundary-regularity condition is imposed on $A$.
\end{abstract}

\maketitle

\section{Introduction}
\label{sec:introduction}

We write
\[
  \D=\{z\in\C:\lvert z\rvert<1\},
  \qquad
  \Torus=\partial\D,
\]
use $X^*$ for the adjoint (the conjugate transpose in finite
dimensions), put $\RePart X=(X+X^*)/2$, and let $\alg(T)$ denote the
unital algebra generated by $T$.

For $n\geq1$, let $\Mn$ be the algebra of $n\times n$ complex
matrices.  For $A\in\Mn$, its numerical range is
\begin{equation}
  W(A)=\{x^*Ax:x\in\C^n,\ \lVert x\rVert_2=1\}.
  \label{eq:numerical-range}
\end{equation}
Throughout, $\lVert x\rVert_2$ denotes the Euclidean norm of a
coordinate vector, whereas $\lVert X\rVert$ denotes the induced
operator norm of a matrix or bounded operator.  Scalar moduli are
written $\lvert\,\cdot\,\rvert$.
The Toeplitz--Hausdorff theorem says that $W(A)$ is convex; in finite
dimensions it is also compact and contains $\sigma(A)$.  The numerical
range therefore records both spectral and nonnormal information.  The
problem considered here is whether it controls the full holomorphic
functional calculus with a universal, sharp constant.

A compact set $K\subset\C$ containing $\sigma(A)$ is called a
$C$-spectral set for $A$ if
\begin{equation}
  \lVert r(A)\rVert\leq C\max_{z\in K}\lvert r(z)\rvert
  \label{eq:spectral-set-definition}
\end{equation}
for every rational function $r$ with no pole on $K$.  Crouzeix proved
the constant-$2$ assertion for $2\times2$ matrices and formulated the
general conjecture in \cite{Crouzeix2004}.  His subsequent
dimension-free theorem gave the constant $11.08$
\cite{Crouzeix2007}.  After further structural and numerical work
\cite{Crouzeix2016}, Crouzeix and Palencia proved that $W(A)$ is a
$(1+\sqrt2)$-spectral set \cite{CrouzeixPalencia2017}.  Their
double-layer and Cauchy-transform
argument was clarified in \cite{RansfordSchwenninger2018} and developed
in several directions in \cite{CaldwellGreenbaumLi2018,CrouzeixGreenbaum2019}.

The intervening literature established the conjectured constant for a
variety of special classes, among them perturbed Jordan blocks, certain
$3\times3$ nilpotent matrices, and tridiagonal $3\times3$ matrices with
elliptic numerical range
\cite{ChoiGreenbaum2012,Crouzeix2013,Choi2013,GladerKurulaLindstrom2018}.
Numerical and extremal investigations appear in
\cite{GreenbaumOverton2018,Li2020}, and a broad account of the problem
and its variants is given in \cite{BickelEtAl2020}.  More recent work
has treated Blaschke-product level sets, cyclicity reductions,
compressions of shifts, further nilpotent classes, and quantitative
bounds for KMS matrices
\cite{BickelGorkin2024,OLoughlinVirtanen2024,BickelEtAl2024,
CrouzeixGreenbaumLi2024}.  A new proof for a class of weighted shifts
appears in \cite{CrouzeixGreenbaum2025}.  None of these reductions
provided a dimension-free route from $1+\sqrt2$ to $2$.

The main result is the following sharp estimate.

\begin{theorem}[Main theorem]
\label{thm:main}
Let $n\geq1$, let $A\in\Mn$, and let $p\in\C[z]$.  Then
\begin{equation}
  \lVert p(A)\rVert
  \leq2\max_{z\in W(A)}\lvert p(z)\rvert.
  \label{eq:main-bound}
\end{equation}
Consequently, $W(A)$ is a $2$-spectral set for $A$.
\end{theorem}

The statement is unconditional: $A$ and $p$ are arbitrary, and no
diagonalizability, spectral-simplicity, algebra-generation, or boundary
regularity condition is imposed.

The constant cannot be decreased.  Indeed, for
\begin{equation}
  N=\begin{pmatrix}0&2\\0&0\end{pmatrix}
  \quad\text{and}\quad p(z)=z,
  \label{eq:sharp-example}
\end{equation}
one has $W(N)=\overline\D$ and $\lVert N\rVert=2$.

\subsection{Where the earlier mechanism loses sharpness}

The positive boundary measure used below originates in the
convex-domain integral calculus of Delyon--Delyon
\cite{DelyonDelyon1999}; see also
\cite{BadeaCrouzeixDelyon2006,CrouzeixPalencia2017}.  If $B$ has
numerical range in a smooth convex domain $\Omega$ and $f$ is bounded
by one on $\Omega$, this calculus produces a unital positive map
$\Phi$ and a holomorphic Cauchy companion $g_f$ such that
\begin{equation}
  2\Phi(f)=f(B)+g_f(B)^*.
  \label{eq:intro-symmetrized}
\end{equation}
Thus positivity directly controls a symmetrized functional calculus,
not $f(B)$ by itself.  The Crouzeix--Palencia argument couples the two
terms in \eqref{eq:intro-symmetrized} and gives the sharp conclusion of
its abstract norm mechanism, namely $1+\sqrt2$.  The examples in
\cite{RansfordSchwenninger2018} show that this abstract mechanism alone
cannot give $2$.  Related uniform-algebra abstractions encounter the
same barrier \cite{ClouatreOstermannRansford2023}.  Configuration-
constant refinements give a strict improvement for every fixed
numerical-range shape, but thin quadrilaterals make the relevant
invariant approach its limiting value, so this method yields no
smaller universal constant
\cite{MalmanEtAl2025,SchwenningerDeVries2025}.  Separately, continuity
and compactness give a constant $C_n<1+\sqrt2$ in every fixed matrix
dimension $n$ \cite{MalmanRatio2024}.

The missing information is algebraic.  The companion in
\eqref{eq:intro-symmetrized} is not an arbitrary operator: for the
approximants constructed in Section~\ref{sec:passage}, it belongs to
the algebra generated by $f(B)^*$.  Estimating it as an unrelated
summand discards precisely this coupling.

\subsection{The new de-symmetrization mechanism}

Our argument retains the whole Cayley family, rather than the single
identity \eqref{eq:intro-symmetrized}.  It has four steps.

\begin{enumerate}
\item Applying $\Phi$ to
\[
  \zeta\longmapsto\frac{1+w f(\zeta)}{1-w f(\zeta)},
  \qquad w\in\D,
\]
produces a matrix-valued Carath\'eodory function $H$, meaning that
$H(0)=I$ and $\RePart H(w)\succeq0$.

\item At an auxiliary stage, a vanishing perturbation arranges that
$T=f(B)$ and $\alg(T)=\alg(B)$.  The double-layer identity then gives
the positive-real completion
\begin{equation}
  H(w)-(I-wT)^{-1}\in\alg(T^*),\qquad w\in\D.
  \label{eq:intro-completion}
\end{equation}

\item For the simple-spectrum approximants used at that stage,
diagonalization turns the unknown term in
\eqref{eq:intro-completion} into a diagonal analytic correction.  We
sample the Herglotz kernel of $H$ at half the conjugate eigenvalues and
add one vector sample at the origin.  The latter is chosen to cancel
the diagonal correction exactly.

\item The surviving positive-matrix inequality compares two weighted
Gramians.  Its difference is positive, and an eigenvector argument
followed by a Stein identity yields $\lVert T\rVert\leq2$.
\end{enumerate}

A perturbation of $f$ that vanishes in the limit enforces algebra
generation on each simple-spectrum approximant.  Successive limits in
the function, the matrix, and the surrounding domain then yield the
estimate for the originally arbitrary $A$ and $p$; no auxiliary
condition remains in the conclusion.

\section{Positive-real kernels}
\label{sec:herglotz}

All matrix inequalities are in the Hermitian order.

A kernel $\mathcal K:\D\times\D\to\Mn$ is positive if, for every
finite choice $w_0,\ldots,w_m\in\D$ and
$\xi_0,\ldots,\xi_m\in\C^n$,
\begin{equation}
  \sum_{i,j=0}^m
  \xi_i^*\mathcal K(w_i,w_j)\xi_j\geq0.
  \label{eq:positive-kernel-definition}
\end{equation}

We shall use the following standard matrix-valued form of the Herglotz
theorem.  Its kernel consequence is included to fix every normalization.

\begin{lemma}[Matrix Herglotz kernel]
\label{lem:herglotz-kernel}
Let $F:\D\to\Mn$ be analytic with
$\RePart F(w)\succeq0$ for every $w\in\D$.  Then
\begin{equation}
  \mathcal L_F(w,z)
  =\frac{F(w)+F(z)^*}{1-w\overline z}
  \label{eq:herglotz-kernel}
\end{equation}
is a positive kernel.
\end{lemma}

\begin{proof}
The matrix Herglotz representation gives a positive
$\Mn$-valued measure $M$ on $\Torus$, and a Hermitian matrix $J_F$, such
that
\begin{equation}
  F(w)=iJ_F+\int_{\Torus}\frac{\zeta+w}{\zeta-w}\,dM(\zeta).
  \label{eq:matrix-herglotz-representation}
\end{equation}
For completeness, this follows by applying the scalar Herglotz theorem
to $u^*F(\,\cdot\,)u$ for every $u\in\C^n$ and polarizing the resulting
positive measures.  Since the space is finite-dimensional, the
polarized scalar measures are the entries of a unique finite
matrix-valued measure; positivity follows from its quadratic forms.

For $\lvert\zeta\rvert=1$, direct calculation gives
\begin{equation}
  \frac{1}{1-w\overline z}\left(
   \frac{\zeta+w}{\zeta-w}
   +\overline{\frac{\zeta+z}{\zeta-z}}
  \right)
  =\frac{2}{(1-w\overline\zeta)(1-\overline z\zeta)}.
  \label{eq:scalar-herglotz-factorization}
\end{equation}
The constant $iJ_F$ cancels from \eqref{eq:herglotz-kernel}.  Substituting
\eqref{eq:matrix-herglotz-representation} and then summing as in
\eqref{eq:positive-kernel-definition} yields
\[
  2\int_{\Torus}
  \left(\sum_i\frac{\xi_i}{1-\overline{w_i}\zeta}\right)^*
  dM(\zeta)
  \left(\sum_j\frac{\xi_j}{1-\overline{w_j}\zeta}\right)
  \geq0.
\]
This proves the assertion.
\end{proof}

\section{The positive-real completion principle}
\label{sec:completion}

The next auxiliary theorem is the new core of the proof.  It converts
an algebraically constrained positive-real completion of a resolvent
into the sharp norm bound.  Section~\ref{sec:double-layer} constructs
the required function at the algebra-generating stage, and
Section~\ref{sec:passage} reaches that stage by perturbation before
passing to arbitrary matrices and polynomials.

\begin{theorem}[Positive-real completion]
\label{thm:completion}
Let $T\in\Mn$ have $n$ distinct eigenvalues, all in $\overline\D$.
Let $H:\D\to\Mn$ be analytic and satisfy
\begin{equation}
  H(0)=I,\qquad \RePart H(w)\succeq0\quad(w\in\D),
  \label{eq:completion-positive-real}
\end{equation}
and
\begin{equation}
  H(w)-(I-wT)^{-1}\in\alg(T^*)\qquad(w\in\D).
  \label{eq:completion-algebra-condition}
\end{equation}
Then $\lVert T\rVert\leq2$.
\end{theorem}

\begin{proof}
Choose a diagonalization
\begin{equation}
  T=S\Lambda S^{-1},\qquad
  \Lambda=\diag(\lambda_1,\ldots,\lambda_n),\qquad
  G=S^*S.
  \label{eq:T-diagonalization}
\end{equation}
The eigenvalues $\lambda_j$ are distinct and belong to
$\overline\D$.  Polynomial interpolation shows that
\begin{equation}
  \alg(T^*)
  =\{(S^{-1})^*DS^*:D\text{ is diagonal}\}.
  \label{eq:adjoint-algebra-diagonal}
\end{equation}
The inverse of the displayed linear parametrization is fixed and
linear.  Hence \eqref{eq:completion-algebra-condition}, applied to an
analytic matrix function, gives an analytic diagonal function
\[
  \Theta(w)=\diag(\theta_1(w),\ldots,\theta_n(w)),
  \qquad \Theta(0)=0,
\]
such that, on setting $\widetilde H(w)=S^*H(w)S$,
\begin{equation}
  \widetilde H(w)=G(I-w\Lambda)^{-1}+\Theta(w)G.
  \label{eq:Htilde-decomposition}
\end{equation}
Congruence by $S$ preserves positive real parts.  Consequently,
Lemma~\ref{lem:herglotz-kernel} applies to $\widetilde H$.

Define Hermitian matrices $P,Q,Y\in\Mn$ by
\begin{equation}
  P_{ij}=\frac{G_{ij}}{1-\overline{\lambda_i}\lambda_j/4},
  \qquad
  Q_{ij}=\frac{G_{ij}}{1-\overline{\lambda_i}\lambda_j/2},
  \qquad
  Y=Q-P.
  \label{eq:PQY-definition}
\end{equation}
Let $e_1,\ldots,e_n$ be the standard basis of $\C^n$, and fix
$u=(u_1,\ldots,u_n)^T\in\C^n$.  In the positive-kernel inequality use
the $n$ points and vectors
\begin{equation}
  w_i=\frac{\overline{\lambda_i}}2,\qquad
  \xi_i=u_i e_i\qquad(1\leq i\leq n),
  \label{eq:eigenvalue-samples}
\end{equation}
and add the point $w_0=0$ with vector
\begin{equation}
  \xi_0=v=-G^{-1}Pu.
  \label{eq:compensating-vector}
\end{equation}
We first verify that this sample cancels the entire unknown function
$\Theta$.  The contribution of the second term in
\eqref{eq:Htilde-decomposition} to the kernel quadratic form is
\begin{equation}
  2\RePart\sum_{i=1}^n
  \overline{u_i}\theta_i(w_i)
  \left(
    (Gv)_i+
    \sum_{j=1}^n
    \frac{G_{ij}u_j}{1-w_i\overline{w_j}}
  \right).
  \label{eq:theta-contribution}
\end{equation}
Because
\[
  1-w_i\overline{w_j}
  =1-\overline{\lambda_i}\lambda_j/4,
\]
the parenthesis in \eqref{eq:theta-contribution} is
$(Gv+Pu)_i$, which vanishes by
\eqref{eq:compensating-vector}.

It remains to calculate the contribution of
$G(I-w\Lambda)^{-1}$.  Between $w_i$ and $w_j$ its $(i,j)$ entry is
\begin{equation}
  \frac{2G_{ij}}
  {(1-\overline{\lambda_i}\lambda_j/2)
   (1-\overline{\lambda_i}\lambda_j/4)}
  =(4Q-2P)_{ij}.
  \label{eq:sampled-main-block}
\end{equation}
Between $w_i$ and $0$ the corresponding row is that of $G+Q$, and
the block at $(0,0)$ is $2G$.  Kernel positivity and
\eqref{eq:compensating-vector} therefore imply
\begin{equation}
\begin{split}
 0\leq{}&u^*(4Q-2P)u
  +2\RePart\bigl(u^*(G+Q)v\bigr)+2v^*Gv\\
 ={}&u^*\bigl(4Y-YG^{-1}P-PG^{-1}Y\bigr)u.
\end{split}
\label{eq:pre-gramian-inequality}
\end{equation}
As $u$ was arbitrary,
\begin{equation}
  4Y-YG^{-1}P-PG^{-1}Y\succeq0.
  \label{eq:Y-inequality}
\end{equation}

Balance the diagonalization by defining
\begin{equation}
\begin{aligned}
  \widetilde T&=G^{1/2}\Lambda G^{-1/2},
  &\widehat P&=G^{-1/2}PG^{-1/2},\\
  \widehat Q&=G^{-1/2}QG^{-1/2},
  &\widehat Y&=\widehat Q-\widehat P.
\end{aligned}
\label{eq:balanced-definitions}
\end{equation}
Expanding the scalar kernels in \eqref{eq:PQY-definition} into
geometric series gives
\begin{equation}
  \widehat P=\sum_{k=0}^{\infty}
  4^{-k}\widetilde T^{*k}\widetilde T^k,
  \qquad
  \widehat Q=\sum_{k=0}^{\infty}
  2^{-k}\widetilde T^{*k}\widetilde T^k.
  \label{eq:two-gramians}
\end{equation}
Indeed, $\widetilde T^k=G^{1/2}\Lambda^kG^{-1/2}$, so the sequence
$(\widetilde T^k)$ is bounded even when some
$\lvert\lambda_j\rvert=1$; the
geometric weights give norm convergence.  It follows that
\begin{equation}
  \widehat Y=
  \sum_{k=1}^{\infty}(2^{-k}-4^{-k})
  \widetilde T^{*k}\widetilde T^k\succeq0.
  \label{eq:Yhat-positive}
\end{equation}
Congruencing \eqref{eq:Y-inequality} by $G^{-1/2}$ yields
\begin{equation}
  4\widehat Y-\widehat Y\widehat P
  -\widehat P\widehat Y\succeq0.
  \label{eq:Yhat-inequality}
\end{equation}

Let $e$ be an eigenvector of the Hermitian matrix $\widehat P$, say
$\widehat Pe=\alpha e$.  If $\alpha>2$, then
\eqref{eq:Yhat-inequality} gives
\begin{equation}
  0\leq2(2-\alpha)e^*\widehat Ye.
  \label{eq:eigenvector-evaluation}
\end{equation}
By \eqref{eq:Yhat-positive}, this forces
$e^*\widehat Ye=0$.  Positivity then gives $\widehat Ye=0$; in
particular, the $k=1$ summand in \eqref{eq:Yhat-positive} gives
$\widetilde Te=0$.  Formula \eqref{eq:two-gramians} now gives
$\widehat Pe=e$, contradicting $\alpha>2$.  The first series in
\eqref{eq:two-gramians} also begins with $I$, and hence
\begin{equation}
  I\preceq\widehat P\preceq2I.
  \label{eq:Phat-bounds}
\end{equation}

Finally, the first Gramian satisfies the Stein identity
\begin{equation}
  \widehat P
  =I+\frac14\widetilde T^*\widehat P\widetilde T.
  \label{eq:stein-identity}
\end{equation}
Using \eqref{eq:Phat-bounds}, we obtain
\[
  \widetilde T^*\widetilde T
  \preceq \widetilde T^*\widehat P\widetilde T
  =4(\widehat P-I)\preceq4I.
\]
Thus $\lVert\widetilde T\rVert\leq2$.  The polar decomposition
$S=UG^{1/2}$ gives
\[
  T=U\widetilde TU^*,
\]
where $U$ is unitary.  Therefore $\lVert T\rVert\leq2$.
\end{proof}

\begin{remark}
\label{rem:completion-novelty}
The factor $1/2$ in the sampling points
\eqref{eq:eigenvalue-samples} simultaneously produces the weights
$1/2$ and $1/4$ in \eqref{eq:two-gramians}.  Their positive difference
and the Stein identity are what recover the sharp constant.  The
sample at the origin is not an estimate: it removes the unknown
analytic correction exactly.
\end{remark}

\section{The double-layer completion}
\label{sec:double-layer}

We now construct from the numerical range the function required by
Theorem~\ref{thm:completion}.  The positive measure and the companion
transform in this section are classical
\cite{DelyonDelyon1999,Crouzeix2007,CrouzeixPalencia2017}; the use of
their full Cayley family to obtain
\eqref{eq:completion-algebra-condition} is the point needed here.

\begin{lemma}[Double-layer positive map]
\label{lem:double-layer-map}
Let $B\in\Mn$, and let $\Omega\subset\C$ be a bounded convex domain
with continuously differentiable boundary and
$W(B)\subset\Omega$.  There is a unital positive map
\[
  \Phi:C(\partial\Omega)\longrightarrow\Mn
\]
such that, whenever $h$ is holomorphic on a neighborhood of
$\overline\Omega$,
\begin{equation}
  2\Phi(h)=h(B)+g_h(B)^*,
  \label{eq:double-layer-identity}
\end{equation}
where
\begin{equation}
  g_h(z)=\frac{1}{2\pi i}
  \int_{\partial\Omega}
  \frac{\overline{h(\zeta)}}{\zeta-z}\,d\zeta,
  \qquad z\in\Omega.
  \label{eq:companion-transform}
\end{equation}
\end{lemma}

\begin{proof}
Orient $\partial\Omega$ counterclockwise.  Let $\nu(\zeta)$ be the
outward unit normal, let $ds$ be arclength, and put
\[
  Z_\zeta=\zeta I-B,
  \qquad
  \mathcal R_\zeta=Z_\zeta^{-1}.
\]
This resolvent exists because
$\sigma(B)\subset W(B)\subset\Omega$.  Define the matrix-valued
boundary measure
\begin{equation}
  d\mu_B(\zeta)=\frac{1}{2\pi}
  \left(
    \nu(\zeta)\mathcal R_\zeta+
    \overline{\nu(\zeta)}\mathcal R_\zeta^*
  \right)ds.
  \label{eq:double-layer-measure}
\end{equation}
It is positive.  In fact, the supporting-line property of the convex
set $\overline\Omega$ gives
\begin{equation}
  \nu(\zeta)Z_\zeta^*+\overline{\nu(\zeta)}Z_\zeta
  =2\RePart\bigl(\overline{\nu(\zeta)}Z_\zeta\bigr)\succeq0.
  \label{eq:supporting-half-plane}
\end{equation}
Multiplication on the left by $\mathcal R_\zeta^*$ and on the right by
$\mathcal R_\zeta$ turns the left side of
\eqref{eq:supporting-half-plane} into the matrix in parentheses in
\eqref{eq:double-layer-measure}.

Along the counterclockwise boundary,
$d\zeta=i\nu(\zeta)ds$.  The matrix Cauchy formula gives
\begin{equation}
  \frac{1}{2\pi}\int_{\partial\Omega}
  \nu(\zeta)\mathcal R_\zeta\,ds=I.
  \label{eq:first-layer-mass}
\end{equation}
Taking adjoints gives the same identity for the second layer.  Hence
\begin{equation}
  \int_{\partial\Omega}d\mu_B=2I.
  \label{eq:double-layer-mass}
\end{equation}
Consequently,
\begin{equation}
  \Phi(\varphi)=\frac12
  \int_{\partial\Omega}\varphi(\zeta)\,d\mu_B(\zeta)
  \label{eq:Phi-definition}
\end{equation}
is unital and positive.  It is therefore star preserving.

For holomorphic $h$, the first layer of $2\Phi(h)$ is $h(B)$ by
Cauchy's formula.  The function $g_h$ in
\eqref{eq:companion-transform} is holomorphic on $\Omega$, so
$g_h(B)$ is defined by the holomorphic functional calculus.  To verify
the relevant integral formula without assuming that $g_h$ extends
across $\partial\Omega$, choose a positively oriented contour
$\Gamma\subset\Omega$ enclosing $\sigma(B)$.  Cauchy's formula and
Fubini's theorem give
\begin{equation}
\begin{split}
 g_h(B)
 &=\frac{1}{2\pi i}\int_{\Gamma}
   g_h(z)(zI-B)^{-1}\,dz\\
 &=\frac{1}{2\pi i}\int_{\partial\Omega}
   \overline{h(\zeta)}(\zeta I-B)^{-1}\,d\zeta.
\end{split}
\label{eq:companion-functional-calculus}
\end{equation}
Taking adjoints in \eqref{eq:companion-functional-calculus} and using
$d\overline\zeta=-i\overline{\nu(\zeta)}ds$ gives the second layer of
$2\Phi(h)$.  This proves \eqref{eq:double-layer-identity}.
\end{proof}

The next construction is used only for auxiliary functions satisfying
the displayed algebra identity.  Section~\ref{sec:passage} enforces
that identity by a vanishing perturbation and then passes to the
original polynomial.

\begin{proposition}[Cayley completion for an auxiliary pair]
\label{prop:cayley-completion}
Let $B$ and $\Omega$ be as in Lemma~\ref{lem:double-layer-map}, and let $f$ be
holomorphic on a neighborhood of $\overline\Omega$ and satisfy
\begin{equation}
  \max_{\zeta\in\overline\Omega}\lvert f(\zeta)\rvert\leq1.
  \label{eq:f-bounded-by-one}
\end{equation}
Put $T=f(B)$.  If
\begin{equation}
  \alg(T)=\alg(B),
  \label{eq:algebra-generation}
\end{equation}
then there is an analytic $H:\D\to\Mn$ satisfying
\begin{equation}
  H(0)=I,\qquad \RePart H(w)\succeq0\quad(w\in\D),
  \label{eq:cayley-H-positive}
\end{equation}
and
\begin{equation}
  H(w)-(I-wT)^{-1}\in\alg(T^*)\qquad(w\in\D).
  \label{eq:cayley-H-completion}
\end{equation}
Moreover, $\sigma(T)\subset\overline\D$.
\end{proposition}

\begin{proof}
Spectral inclusion follows from
$\sigma(B)\subset\Omega$, spectral mapping, and
\eqref{eq:f-bounded-by-one}.  For $w\in\D$, define on the boundary
\begin{equation}
  c_w(\zeta)=\frac{1+w f(\zeta)}{1-w f(\zeta)}
  \label{eq:cayley-boundary-function}
\end{equation}
and put
\begin{equation}
  H(w)=\Phi(c_w).
  \label{eq:H-from-Phi}
\end{equation}
The map $w\mapsto c_w$ is analytic with values in
$C(\partial\Omega)$, locally uniformly on $\D$.  Hence $H$ is
analytic, and $H(0)=I$.  Furthermore,
\[
  \RePart c_w(\zeta)
  =\frac{1-\lvert w f(\zeta)\rvert^2}
  {\lvert1-w f(\zeta)\rvert^2}\geq0.
\]
The positivity and star preservation of $\Phi$ give
$\RePart H(w)=\Phi(\RePart c_w)\succeq0$.

The uniformly convergent boundary expansion
\begin{equation}
  c_w=1+2\sum_{m=1}^{\infty}w^m f^m
  \label{eq:cayley-series}
\end{equation}
and \eqref{eq:double-layer-identity} yield, with norm convergence,
\begin{equation}
  H(w)=I+\sum_{m=1}^{\infty}w^m
  \bigl(T^m+g_{f^m}(B)^*\bigr).
  \label{eq:H-series}
\end{equation}
Since the spectral radius of $wT$ is less than $1$, the Neumann series
for $(I-wT)^{-1}$ converges in norm.  Subtracting it from
\eqref{eq:H-series} gives
\begin{equation}
  H(w)-(I-wT)^{-1}
  =\sum_{m=1}^{\infty}w^m g_{f^m}(B)^*.
  \label{eq:completion-series}
\end{equation}
Each $g_{f^m}$ is holomorphic on a neighborhood of $\sigma(B)$.
Finite-dimensional holomorphic functional calculus, equivalently
Hermite interpolation followed by Cayley--Hamilton, implies
$g_{f^m}(B)\in\alg(B)$.  By
\eqref{eq:algebra-generation}, every partial sum on the right of
\eqref{eq:completion-series} belongs to $\alg(T^*)$.  This algebra is
finite-dimensional and therefore norm closed, which proves
\eqref{eq:cayley-H-completion}.
\end{proof}

Combining the preceding results gives the sharp estimate at the
auxiliary simple-spectrum, algebra-generating stage.  The next section
constructs this stage arbitrarily close to the original data and then
passes to the limit.

\begin{corollary}[Auxiliary sharp bound]
\label{cor:auxiliary-bound}
In Proposition~\ref{prop:cayley-completion}, if $T=f(B)$ has distinct
eigenvalues, then $\lVert f(B)\rVert\leq2$.
\end{corollary}

\begin{proof}
Apply Theorem~\ref{thm:completion} to the function $H$ constructed in
Proposition~\ref{prop:cayley-completion}.
\end{proof}

\section{Passage to arbitrary matrices and proof of the main theorem}
\label{sec:passage}

We now pass from the auxiliary stage to the arbitrary matrix and
polynomial in Theorem~\ref{thm:main}.  Two elementary approximation
facts supply simple-spectrum matrices and smooth surrounding domains;
the proof then enforces algebra generation by a vanishing perturbation.

\begin{lemma}[Simple spectrum and numerical range]
\label{lem:simple-spectrum-density}
Matrices with $n$ distinct eigenvalues are dense in $\Mn$.  Moreover,
for $A,B\in\Mn$,
\begin{equation}
  W(B)\subset
  W(A)+\{z\in\C:\lvert z\rvert\leq\lVert B-A\rVert\}.
  \label{eq:numerical-range-Lipschitz}
\end{equation}
\end{lemma}

\begin{proof}
The discriminant of the characteristic polynomial is a nonzero
polynomial in the matrix entries.  Its nonvanishing set is therefore
dense in $\Mn$, proving the first assertion.  For a unit vector $x$,
\[
  \lvert x^*(B-A)x\rvert\leq\lVert B-A\rVert,
\]
which proves \eqref{eq:numerical-range-Lipschitz}.
\end{proof}

\begin{lemma}[Outer approximation]
\label{lem:smooth-outer-approximation}
If $K\subset\C$ is nonempty, compact, and convex, then there are
bounded convex domains $\Omega_j$ with continuously differentiable
boundaries such that
\begin{equation}
  K\subset\Omega_j,
  \qquad
  d_{\mathrm H}(\overline\Omega_j,K)\longrightarrow0,
  \label{eq:outer-Hausdorff}
\end{equation}
where $d_{\mathrm H}$ denotes the Hausdorff distance between compact
sets.
\end{lemma}

\begin{proof}
For $\delta>0$, let
\[
  \Omega_\delta=\{z\in\C:\dist(z,K)<\delta\}.
\]
This is a bounded convex domain and its closure is the Minkowski sum
$K+\delta\overline\D$, so its Hausdorff distance from $K$ is
$\delta$.  The metric projection $\pi_K:\C\to K$ is unique and
continuous.  The squared distance to $K$ is continuously
differentiable off $K$, with gradient
$2(z-\pi_Kz)$.  Its gradient does not vanish on the level set
$\dist(z,K)=\delta$.  The implicit-function theorem therefore shows
that $\partial\Omega_\delta$ is continuously differentiable.  Taking
$\delta\downarrow0$ proves the lemma.
\end{proof}

We now prove Theorem~\ref{thm:main}.

\begin{proof}[Proof of Theorem~\ref{thm:main}]
Let $A$ and $p$ be arbitrary as in Theorem~\ref{thm:main}.  Choose a
bounded convex domain $\Omega$ with continuously differentiable boundary
such that
\begin{equation}
  W(A)\subset\Omega.
  \label{eq:WA-in-Omega}
\end{equation}
We first prove
\begin{equation}
  \lVert p(A)\rVert
  \leq2\max_{z\in\overline\Omega}\lvert p(z)\rvert.
  \label{eq:fixed-Omega-bound}
\end{equation}
Because $W(A)$ is compactly contained in $\Omega$,
\[
  \dist(W(A),\C\setminus\Omega)>0.
\]
Lemma~\ref{lem:simple-spectrum-density} and
\eqref{eq:numerical-range-Lipschitz} therefore give a sequence of
simple-spectrum matrices $B_k\to A$ with
$W(B_k)\subset\Omega$.  Since $p(B_k)\to p(A)$, it is enough to prove
\eqref{eq:fixed-Omega-bound} for an arbitrary such matrix $B$.

Write the distinct eigenvalues of $B$ as
$\beta_1,\ldots,\beta_n$ and set
\begin{equation}
  m_\Omega=\max_{z\in\overline\Omega}\lvert p(z)\rvert.
  \label{eq:m-Omega}
\end{equation}
If $m_\Omega=0$, the polynomial vanishes on the open set $\Omega$ and
is identically zero.  In the remaining case $m_\Omega>0$, put
\begin{equation}
  f=\frac{p}{m_\Omega},
  \qquad
  R_\Omega=\max_{z\in\overline\Omega}\lvert z\rvert.
  \label{eq:normalized-f}
\end{equation}
For $\eta\in\C$ define
\begin{equation}
  f_\eta(z)=
  \frac{f(z)+\eta z}{1+\lvert\eta\rvert R_\Omega}.
  \label{eq:f-eta}
\end{equation}
Then
$\max_{z\in\overline\Omega}\lvert f_\eta(z)\rvert\leq1$.
For a pair $i\neq j$,
the equality
\[
  f(\beta_i)+\eta\beta_i
  =f(\beta_j)+\eta\beta_j
\]
excludes at most one value of $\eta$.  There are thus only finitely
many exceptional values for which the numbers
$f_\eta(\beta_1),\ldots,f_\eta(\beta_n)$ fail to be distinct.  Choose
nonexceptional $\eta_\ell\to0$ and put
\begin{equation}
  T_\ell=f_{\eta_\ell}(B).
  \label{eq:T-ell}
\end{equation}
Each $T_\ell$ has distinct eigenvalues.  It also generates the same
unital algebra as $B$.  Indeed, $T_\ell\in\alg(B)$, while Lagrange
interpolation at the distinct points
$f_{\eta_\ell}(\beta_i)$ supplies a polynomial $q_\ell$ satisfying
\[
  q_\ell(f_{\eta_\ell}(\beta_i))=\beta_i
  \qquad(1\leq i\leq n).
\]
Since $B$ is diagonalizable, this implies
$q_\ell(T_\ell)=B$.  Consequently,
\begin{equation}
  \alg(T_\ell)=\alg(B).
  \label{eq:T-generates-B}
\end{equation}

Corollary~\ref{cor:auxiliary-bound}, applied to $B$ and
$f_{\eta_\ell}$, gives
\[
  \lVert f_{\eta_\ell}(B)\rVert\leq2.
\]
Letting $\ell\to\infty$ yields $\lVert f(B)\rVert\leq2$, or
\begin{equation}
  \lVert p(B)\rVert\leq2m_\Omega.
  \label{eq:B-Omega-bound}
\end{equation}
Letting the simple-spectrum matrices $B_k$ tend to $A$ proves
\eqref{eq:fixed-Omega-bound}.

Apply Lemma~\ref{lem:smooth-outer-approximation} to the compact convex
set $K=W(A)$, obtaining $\Omega_j$.  All these closures eventually
lie in one fixed compact neighborhood of $K$.  Uniform continuity of
$p$ there and \eqref{eq:outer-Hausdorff} give
\begin{equation}
  \max_{z\in\overline\Omega_j}\lvert p(z)\rvert
  \longrightarrow
  \max_{z\in W(A)}\lvert p(z)\rvert.
  \label{eq:maxima-converge}
\end{equation}
Apply \eqref{eq:fixed-Omega-bound} to each $\Omega_j$ and pass to the
limit.  This proves \eqref{eq:main-bound}.

Thus the polynomial estimate holds for the originally arbitrary $A$
and $p$.  The simple-spectrum matrices, the algebra-generating
perturbations, and the smooth surrounding domains occurred only inside
approximating sequences and have all disappeared in the successive
limits.

It remains to justify the spectral-set conclusion.  Let $r$ be
rational with no pole on $W(A)$.  Choose $\delta>0$ such that $r$ is
holomorphic on the closed $\delta$-neighborhood $K_\delta$ of
$W(A)$.  The compact convex set $K_\delta$ has connected complement,
so polynomial Runge approximation gives polynomials $p_j$ converging
uniformly to $r$ on $K_\delta$.  As
$\sigma(A)\subset W(A)\subset\operatorname{int}K_\delta$, the
holomorphic functional calculus gives $p_j(A)\to r(A)$.  Also
\[
  \max_{z\in W(A)}\lvert p_j(z)\rvert
  \longrightarrow
  \max_{z\in W(A)}\lvert r(z)\rvert.
\]
Apply the polynomial estimate to $p_j$ and pass to the limit to obtain
\eqref{eq:spectral-set-definition} with $K=W(A)$ and $C=2$.
\end{proof}

\section{Consequences and scope}
\label{sec:consequences}

Theorem~\ref{thm:main} immediately gives the usual holomorphic version
on a neighborhood of the numerical range.

\begin{corollary}
\label{cor:holomorphic-bound}
Let $A\in\Mn$, and let $f$ be holomorphic on a neighborhood of
$W(A)$.  Then
\begin{equation}
  \lVert f(A)\rVert
  \leq2\max_{z\in W(A)}\lvert f(z)\rvert.
  \label{eq:holomorphic-bound}
\end{equation}
\end{corollary}

\begin{proof}
Polynomial Runge approximation on a sufficiently small closed convex
neighborhood of $W(A)$ gives the same limiting argument used at the
end of the proof of Theorem~\ref{thm:main}.
\end{proof}

The operator form follows by the standard finite-dimensional
compression argument; compare \cite{BickelEtAl2020}.

\begin{corollary}[Hilbert-space form]
\label{cor:operator-bound}
Let $\mathcal H$ be a complex Hilbert space and let
$A\in\mathcal B(\mathcal H)$.  For every $p\in\C[z]$,
\begin{equation}
  \lVert p(A)\rVert
  \leq 2\sup_{z\in W(A)}\lvert p(z)\rvert.
  \label{eq:operator-polynomial-bound}
\end{equation}
Consequently, $\overline{W(A)}$ is a $2$-spectral set for $A$.
\end{corollary}

\begin{proof}
The assertion is immediate if $\mathcal H=\{0\}$ or $p=0$.  Otherwise,
write $d=\deg p$.  For a unit vector $x\in\mathcal H$, let
\[
  \mathcal M_x=\operatorname{span}\{x,Ax,\ldots,A^d x\},
  \qquad
  A_x=\Pi_x\bigl(A\!\restriction_{\mathcal M_x}\bigr),
\]
where $\Pi_x$ is the orthogonal projection onto
$\mathcal M_x$.  Induction gives
\begin{equation}
  A_x^k x=A^k x \qquad(0\leq k\leq d),
  \label{eq:Krylov-agreement}
\end{equation}
because $A^{k+1}x\in\mathcal M_x$ whenever $k<d$.  Moreover,
$W(A_x)\subset W(A)$.  Theorem~\ref{thm:main}, applied on the
finite-dimensional space $\mathcal M_x$, therefore gives
\[
  \lVert p(A)x\rVert_{\mathcal H}
  =\lVert p(A_x)x\rVert_{\mathcal H}
  \leq\lVert p(A_x)\rVert
  \leq2\sup_{z\in W(A)}\lvert p(z)\rvert.
\]
Taking the supremum over unit vectors $x$ proves
\eqref{eq:operator-polynomial-bound}.

Put $K=\overline{W(A)}$.  The standard inclusion
$\sigma(A)\subset K$ shows that the rational functional calculus is
defined for every rational function $r$ without poles on $K$.  Choose
$\delta>0$ so that $r$ is holomorphic on the closed
$\delta$-neighborhood $K_\delta$ of $K$.  This set is compact and
convex, so polynomial Runge approximation gives polynomials $p_j$
converging uniformly to $r$ on $K_\delta$.  The holomorphic functional
calculus gives $p_j(A)\to r(A)$, while
\[
  \sup_{z\in W(A)}\lvert p_j(z)\rvert
  \longrightarrow
  \sup_{z\in W(A)}\lvert r(z)\rvert.
\]
Passing to the limit in \eqref{eq:operator-polynomial-bound} proves the
rational estimate and hence the spectral-set assertion.
\end{proof}

The conclusion controls errors in polynomial and rational
approximations of matrix functions: if $q$ approximates $f$ uniformly
on $W(A)$, then
\begin{equation}
  \lVert f(A)-q(A)\rVert
  \leq2\max_{z\in W(A)}\lvert f(z)-q(z)\rvert.
  \label{eq:matrix-function-error}
\end{equation}
This is the sharp universal form of the numerical-range estimate that
underlies applications to GMRES and rational Krylov methods; see
\cite{ChoiGreenbaum2015,CrouzeixGreenbaum2019,ChenGreenbaumTrogdon2025}
for the role of spectral sets in those settings.

\section{Discussion and further directions}
\label{sec:outlook}

Theorem~\ref{thm:main} proves exactly the original Crouzeix conjecture.
Beyond this theorem, one may ask for a
\emph{uniform matrix-polynomial extension}: for every
$m\geq1$ and every matrix polynomial
$\mathcal F=[f_{ij}]\in\mathbb M_m(\C[z])$, does one have
\begin{equation}
  \lVert\mathcal F(A)\rVert
  \leq2\max_{z\in W(A)}\lVert\mathcal F(z)\rVert,
  \qquad \mathcal F(A)=[f_{ij}(A)].
  \label{eq:uniform-matrix-polynomial-extension}
\end{equation}
Here $\lVert\mathcal F(A)\rVert$ is the operator norm on $(\C^n)^m$,
whereas $\lVert\mathcal F(z)\rVert$ is the operator norm on $\C^m$.
The case $m=1$ is precisely the original conjecture proved above;
uniformity in $m$ is an additional operator-algebraic requirement.  The
foundational matrix-level disk estimate is due to Okubo and Ando
\cite{OkuboAndo1975}.  Uniform matrix-level spectral-set bounds and
their relation to numerical radius are studied in
\cite{DavidsonPaulsenWoerdeman2018,HartzMcCarthy2026}; the constant
$1+\sqrt2$ at every matrix level follows from
\cite{CrouzeixPalencia2017}, while sharp $\rho$-contraction estimates
are developed in \cite{SchwenningerDeVriesRho2025}.  The extension is
known with constant $2$ for weighted shift matrices
\cite{CrouzeixGreenbaum2025}.  The auxiliary simple-spectrum
interpolation in Section~\ref{sec:passage} is eliminated by the
limiting argument that proves Theorem~\ref{thm:main}.  At the uniform
matrix-polynomial level, however, the corresponding diagonal-algebra
mechanism is not stable in the polynomial matrix size, making this a
natural separate problem in operator-algebraic functional calculus.

The positive-real completion principle is independent of numerical
ranges: it starts with a Carath\'eodory function whose resolvent defect
lies in an adjoint-generated algebra.  This separation suggests
several concrete problems.

\begin{enumerate}
\item \emph{Uniform matrix-polynomial extension.}
Find a version of the diagonal-correction argument that is stable
uniformly in the polynomial matrix size.  Matrix-polynomial inputs lead
to block coefficients that need not reduce to a single diagonal
algebra, so the cancellation in \eqref{eq:theta-contribution} requires
a genuinely operator-algebraic formulation.

\item \emph{Infinite-dimensional completion theory.}
A direct infinite-dimensional version of Theorem~\ref{thm:completion},
based for example on spectral measures or Riesz bases, could reveal
equality cases or uniform
matrix-level functional-calculus structure.  The main obstacle is to
replace the finite diagonal correction and Lagrange generation
\eqref{eq:T-generates-B} by a stable functional-calculus condition
without relying on a diagonal spectral model.
Related questions for unbounded numerical ranges have recently been
studied in \cite{CrouzeixUnbounded2025}.

\item \emph{Abstract uniform algebras.}
The additional unitality structure isolated in
\cite{ClouatreOstermannRansford2023,HartzMcCarthy2026} is close in
spirit to
\eqref{eq:completion-algebra-condition}.  It is natural to ask whether
positive-real completions yield the conjectured sharp constant in that
abstract setting, or sharpen the spectral constants studied in
\cite{SchwenningerDeVries2024,BlazhkoEtAl2025}.

\item \emph{Other geometric functional calculi.}
Double-layer methods apply to domains beyond numerical ranges and have
recently been recast through configuration constants and annular
double-layer kernels
\cite{CrouzeixGreenbaum2019,SchwenningerDeVries2025,MalmanEtAl2025,
JuryTsikalas2024,AglerLykovaYoung2025}.  Elliptic domains also admit
structured dilation models \cite{AglerLykovaYoung2024}.  Whenever the
companion
transform lands in a controlled adjoint algebra, the kernel-sampling
argument may replace a triangle-inequality estimate.  This offers a
route to sharp constants for annular, multiply connected, or
shape-constrained spectral sets.

\item \emph{Equality and stability.}
With the notation of the proof of Theorem~\ref{thm:completion},
if $\lVert\widetilde T\rVert=2$ and $e\in\C^n$ satisfies
$\lVert e\rVert_2=1$ and $\lVert\widetilde Te\rVert_2=2$, then equality
throughout the last
operator-inequality chain in the proof of
Theorem~\ref{thm:completion} forces
\[
  \widehat Pe=2e,
  \qquad
  \widehat P\widetilde Te=\widetilde Te.
\]
Analyzing these vector saturation conditions may classify extremal
pairs $(A,p)$ and quantify stability near the nilpotent example
\eqref{eq:sharp-example}.  Such a
classification must allow the nonuniqueness phenomena found in
\cite{Li2020}.

\item \emph{Generalized numerical ranges.}
Recent work proposes spectral-set problems for $q$-numerical ranges
\cite{OLoughlinRani2026}.  Their double-layer structure makes them a
natural testing ground for positive-real completion, although no
constant-$2$ analogue is asserted here.
\end{enumerate}

\section*{AI-assistance disclosure}

OpenAI ChatGPT contributed the idea of sampling the matrix Herglotz
kernel at half the conjugate eigenvalues and adding the origin sample
that cancels the adjoint-algebra correction; this appears in the proof
of Theorem~\ref{thm:completion}, especially
\eqref{eq:eigenvalue-samples}--\eqref{eq:theta-contribution}.  It also
assisted with exposition, bibliography, and typesetting.  The human
author assumes full responsibility for the correctness of all
mathematical statements and for the integrity and accuracy of all
citations.

\begin{thebibliography}{99}

\bibitem{AglerLykovaYoung2024}
J. Agler, Z. A. Lykova, and N. J. Young,
\textit{On the operators with numerical range in an ellipse},
J. Funct. Anal. \textbf{287} (2024), no.~8, Article 110556, 62 pp.,
\doi{10.1016/j.jfa.2024.110556}.

\bibitem{AglerLykovaYoung2025}
J. Agler, Z. A. Lykova, and N. J. Young,
\textit{Function theory on the annulus in the dp-norm},
Integral Equations Operator Theory \textbf{97} (2025), Article 30,
\doi{10.1007/s00020-025-02814-w}.

\bibitem{BadeaCrouzeixDelyon2006}
C. Badea, M. Crouzeix, and B. Delyon,
\textit{Convex domains and $K$-spectral sets},
Math. Z. \textbf{252} (2006), no.~2, 345--365,
\doi{10.1007/s00209-005-0857-y}.

\bibitem{BickelEtAl2020}
K. Bickel, P. Gorkin, A. Greenbaum, T. Ransford,
F. L. Schwenninger, and E. Wegert,
\textit{Crouzeix's conjecture and related problems},
Comput. Methods Funct. Theory \textbf{20} (2020), no.~3--4, 701--728,
\doi{10.1007/s40315-020-00350-9}.

\bibitem{BickelEtAl2024}
K. Bickel, G. Corbett, A. Glenning, C. Guan, and M. Vollmayr-Lee,
\textit{Crouzeix's conjecture, compressions of shifts, and classes of
nilpotent matrices},
Concrete Operators \textbf{11} (2024), no.~1, Paper No.~20240004,
\doi{10.1515/conop-2024-0004}.

\bibitem{BickelGorkin2024}
K. Bickel and P. Gorkin,
\textit{Blaschke products, level sets, and Crouzeix's conjecture},
J. Anal. Math. \textbf{152} (2024), no.~1, 217--254,
\doi{10.1007/s11854-023-0295-y}.

\bibitem{BlazhkoEtAl2025}
H. Blazhko, D. Homza, F. L. Schwenninger, J. de Vries, and M. Wojtylak,
\textit{The algebraic numerical range as a spectral set in Banach
algebras},
Canad. J. Math., published online 2025,
\doi{10.4153/S0008414X25000124}.

\bibitem{CaldwellGreenbaumLi2018}
T. Caldwell, A. Greenbaum, and K. Li,
\textit{Some extensions of the Crouzeix--Palencia result},
SIAM J. Matrix Anal. Appl. \textbf{39} (2018), no.~2, 769--780,
\doi{10.1137/17M1140832}.

\bibitem{ChenGreenbaumTrogdon2025}
T. Chen, A. Greenbaum, and T. Trogdon,
\textit{GMRES, pseudospectra, and Crouzeix's conjecture for shifted and
scaled Ginibre matrices},
Math. Comp. \textbf{94} (2025), no.~351, 241--261,
\doi{10.1090/mcom/3963}.

\bibitem{Choi2013}
D. Choi,
\textit{A proof of Crouzeix's conjecture for a class of matrices},
Linear Algebra Appl. \textbf{438} (2013), no.~8, 3247--3257,
\doi{10.1016/j.laa.2012.12.045}.

\bibitem{ChoiGreenbaum2012}
A. Greenbaum and D. Choi,
\textit{Crouzeix's conjecture and perturbed Jordan blocks},
Linear Algebra Appl. \textbf{436} (2012), no.~7, 2342--2352,
\doi{10.1016/j.laa.2011.12.005}.

\bibitem{ChoiGreenbaum2015}
D. Choi and A. Greenbaum,
\textit{Roots of matrices in the study of GMRES convergence and
Crouzeix's conjecture},
SIAM J. Matrix Anal. Appl. \textbf{36} (2015), no.~1, 289--301,
\doi{10.1137/140961742}.

\bibitem{ClouatreOstermannRansford2023}
R. Clou\^atre, M. Ostermann, and T. Ransford,
\textit{An abstract approach to the Crouzeix conjecture},
J. Operator Theory \textbf{90} (2023), no.~1, 209--221,
\doi{10.7900/jot.2021nov15.2364}.

\bibitem{Crouzeix2004}
M. Crouzeix,
\textit{Bounds for analytical functions of matrices},
Integral Equations Operator Theory \textbf{48} (2004), no.~4, 461--477,
\doi{10.1007/s00020-002-1188-6}.

\bibitem{Crouzeix2007}
M. Crouzeix,
\textit{Numerical range and functional calculus in Hilbert space},
J. Funct. Anal. \textbf{244} (2007), no.~2, 668--690,
\doi{10.1016/j.jfa.2006.10.013}.

\bibitem{Crouzeix2013}
M. Crouzeix,
\textit{Spectral sets and $3\times3$ nilpotent matrices},
in C. Badea, D. Li, and V. Petkova (eds.),
\textit{Topics in Functional and Harmonic Analysis},
Theta Ser. Adv. Math., vol.~14, Theta, Bucharest, 2013, pp.~27--42.

\bibitem{Crouzeix2016}
M. Crouzeix,
\textit{Some constants related to numerical ranges},
SIAM J. Matrix Anal. Appl. \textbf{37} (2016), no.~1, 420--442,
\doi{10.1137/15M1020411}.

\bibitem{CrouzeixUnbounded2025}
M. Crouzeix,
\textit{Numerical ranges and spectral sets: the unbounded case},
preprint, 2025, arXiv:2509.19792,
\url{https://arxiv.org/abs/2509.19792}.

\bibitem{CrouzeixGreenbaum2019}
M. Crouzeix and A. Greenbaum,
\textit{Spectral sets: numerical range and beyond},
SIAM J. Matrix Anal. Appl. \textbf{40} (2019), no.~3, 1087--1101,
\doi{10.1137/18M1198417}.

\bibitem{CrouzeixGreenbaumLi2024}
M. Crouzeix, A. Greenbaum, and K. Li,
\textit{Numerical bounds on the Crouzeix ratio for a class of matrices},
Calcolo \textbf{61} (2024), no.~2, Paper No.~32, 18 pp.,
\doi{10.1007/s10092-024-00580-6}.

\bibitem{CrouzeixGreenbaum2025}
M. Crouzeix and A. Greenbaum,
\textit{A new proof that the numerical range is a complete $2$-spectral
set for weighted shift matrices},
preprint, 2025, arXiv:2508.12768,
\url{https://arxiv.org/abs/2508.12768}.

\bibitem{CrouzeixPalencia2017}
M. Crouzeix and C. Palencia,
\textit{The numerical range is a $(1+\sqrt2)$-spectral set},
SIAM J. Matrix Anal. Appl. \textbf{38} (2017), no.~2, 649--655,
\doi{10.1137/17M1116672}.

\bibitem{DavidsonPaulsenWoerdeman2018}
K. R. Davidson, V. I. Paulsen, and H. J. Woerdeman,
\textit{Complete spectral sets and numerical range},
Proc. Amer. Math. Soc. \textbf{146} (2018), no.~3, 1189--1195,
\doi{10.1090/proc/13801}.

\bibitem{DelyonDelyon1999}
B. Delyon and F. Delyon,
\textit{Generalization of von Neumann's spectral sets and integral
representation of operators},
Bull. Soc. Math. France \textbf{127} (1999), no.~1, 25--41,
\doi{10.24033/bsmf.2340}.

\bibitem{GladerKurulaLindstrom2018}
C. Glader, M. Kurula, and M. Lindstr\"om,
\textit{Crouzeix's conjecture holds for tridiagonal $3\times3$ matrices
with elliptic numerical range centered at an eigenvalue},
SIAM J. Matrix Anal. Appl. \textbf{39} (2018), no.~1, 346--364,
\doi{10.1137/17M1110663}.

\bibitem{GreenbaumOverton2018}
A. Greenbaum and M. L. Overton,
\textit{Numerical investigation of Crouzeix's conjecture},
Linear Algebra Appl. \textbf{542} (2018), 225--245,
\doi{10.1016/j.laa.2017.04.035}.

\bibitem{HartzMcCarthy2026}
M. Hartz and J. E. McCarthy,
\textit{From Clou\^atre--Ostermann--Ransford to Okubo--Ando},
preprint, 2026, arXiv:2606.02922,
\url{https://arxiv.org/abs/2606.02922}.

\bibitem{JuryTsikalas2024}
M. T. Jury and G. Tsikalas,
\textit{Positivity conditions on the annulus via the double-layer
potential kernel},
Studia Math. \textbf{278} (2024), no.~3, 233--265,
\doi{10.4064/sm231023-2-8}.

\bibitem{Li2020}
K. Li,
\textit{On the uniqueness of functions that maximize the Crouzeix ratio},
Linear Algebra Appl. \textbf{599} (2020), 105--120,
\doi{10.1016/j.laa.2020.03.045}.

\bibitem{MalmanEtAl2025}
B. Malman, J. Mashreghi, R. O'Loughlin, and T. Ransford,
\textit{Double-layer potentials, configuration constants, and
applications to numerical ranges},
Int. Math. Res. Not. IMRN \textbf{2025} (2025), no.~8,
Paper No.~rnaf084, 34 pp.,
\doi{10.1093/imrn/rnaf084}.

\bibitem{MalmanRatio2024}
B. Malman, J. Mashreghi, R. O'Loughlin, and T. Ransford,
\textit{On the Crouzeix ratio for $N\times N$ matrices},
preprint, 2024, arXiv:2409.14127,
\url{https://arxiv.org/abs/2409.14127}.

\bibitem{OkuboAndo1975}
K. Okubo and T. Ando,
\textit{Constants related to operators of class $C_\rho$},
Manuscripta Math. \textbf{16} (1975), no.~4, 385--394,
\doi{10.1007/BF01323467}.

\bibitem{OLoughlinRani2026}
R. O'Loughlin and J. Rani,
\textit{$q$-numerical ranges and spectral sets},
preprint, 2026, arXiv:2603.15536,
\url{https://arxiv.org/abs/2603.15536}.

\bibitem{OLoughlinVirtanen2024}
R. O'Loughlin and J. Virtanen,
\textit{Crouzeix's conjecture for classes of matrices},
Linear Algebra Appl. \textbf{697} (2024), 277--292,
\doi{10.1016/j.laa.2023.12.008}.

\bibitem{RansfordSchwenninger2018}
T. Ransford and F. L. Schwenninger,
\textit{Remarks on the Crouzeix--Palencia proof that the numerical
range is a $(1+\sqrt2)$-spectral set},
SIAM J. Matrix Anal. Appl. \textbf{39} (2018), no.~1, 342--345,
\doi{10.1137/17M1143757}.

\bibitem{SchwenningerDeVries2024}
F. L. Schwenninger and J. de Vries,
\textit{On abstract spectral constants},
in M. Ptak, H. J. Woerdeman, and M. Wojtylak (eds.),
\textit{Operator and Matrix Theory, Function Spaces, and
Applications}, Oper. Theory Adv. Appl., vol.~295, Birkh\"auser/Springer,
Cham, 2024, pp.~335--350,
\doi{10.1007/978-3-031-50613-0_15}.

\bibitem{SchwenningerDeVries2025}
F. L. Schwenninger and J. de Vries,
\textit{The double-layer potential for spectral constants revisited},
Integral Equations Operator Theory \textbf{97} (2025), no.~2,
Paper No.~13, 22 pp.,
\doi{10.1007/s00020-025-02800-2}.

\bibitem{SchwenningerDeVriesRho2025}
F. L. Schwenninger and J. de Vries,
\textit{A sharp bound for the functional calculus of
$\rho$-contractions},
Proc. Amer. Math. Soc. \textbf{153} (2025), no.~11, 4759--4768,
\doi{10.1090/proc/17346}.

\end{thebibliography}

\end{document}
